\chapter{Experimental Validation}
\label{chap:validation}

The previous chapter described what \nidsname{} does on paper. This
chapter describes what it does in practice. The validation drives
the entire attack catalogue through the dispatcher using the
integrated Attack Console, measures alert generation per attack
kind, and reports the latency from packet injection to UI
notification. The methodology is reproducible: the synthetic
traffic generator uses a fixed random seed and the Attack Console
endpoints are scriptable through the REST API, so the same numbers
can be obtained on any machine that runs the project.

\section{Test Environment}
\label{sec:val-environment}

All measurements in this chapter were collected on a single
laptop with the following specification:

\begin{itemize}
  \item Intel Core i7-12700H, 14 cores, 20 threads, 2.3 GHz base.
  \item 16\,GB DDR5 RAM.
  \item Windows 11 Pro 23H2 (build 22631).
  \item Python 3.12.5 in a project-local virtual environment.
  \item Node.js 20.11 LTS for the frontend bundler.
\end{itemize}

\section{Synthetic Validation}
\label{sec:val-synthetic}

\subsection{Setup}
\label{sec:val-synthetic-setup}

The synthetic scenario is the default deployment of the
application. The environment file \file{backend/.env} has
\code{SIMULATE\_TRAFFIC=true}, \code{SIM\_ATTACKS\_ON\_DEMAND\_ONLY=true}
and \code{SIM\_BENIGN\_TRAFFIC=false}. With these defaults the
dispatcher sees no packets at all until the operator clicks an
attack button in the Attack Console.

I drove the test through the \code{POST /api/sim/attack} endpoint
rather than the UI, both for repeatability and so the results
could be scripted. The Python script in
Listing~\ref{lst:val-script} fires each attack type five times
with a one-second pause between bursts and records the number of
alerts each burst produced.

\begin{lstlisting}[style=pyfile,caption={Synthetic validation driver},label={lst:val-script}]
import time, requests

URL = "http://localhost:8000"
LOGIN = requests.post(URL + "/api/auth/login",
                      data={"username": "admin", "password": "admin"})
TOKEN = LOGIN.json()["access_token"]
H = {"Authorization": "Bearer " + TOKEN}

ATTACKS = ["nmap", "ssh_brute", "icmp_flood",
           "sqli", "cmd_inject", "reverse_shell"]

for kind in ATTACKS:
    for _ in range(5):
        r = requests.post(URL + "/api/sim/attack",
                          json={"kind": kind}, headers=H)
        d = r.json()
        print(f"{kind:14s}  pkts={d['packet_count']:3d}  "
              f"alerts={d['alerts_generated']:3d}")
        time.sleep(1.0)
\end{lstlisting}

\subsection{Results}
\label{sec:val-synthetic-results}

Table~\ref{tab:val-synthetic} summarises the outcome.
The packet counts reflect the simulator's randomised burst sizes,
and the alert counts reflect how many distinct alerts each burst
produced.

\begin{table}[h]
\centering
\caption{Synthetic attack validation: packets injected and alerts
generated, averaged over five runs per attack kind.}
\label{tab:val-synthetic}
\small
\begin{tabular}{l r r r r}
\toprule
Attack kind     & Avg packets & Min alerts & Max alerts & Avg alerts \\
\midrule
nmap            & 32 & 23 & 32 & 27.4 \\
ssh\_brute      & 16 &  9 & 13 & 11.0 \\
icmp\_flood     & 73 & 47 & 65 & 54.6 \\
sqli            &  1 &  1 &  1 &  1.0 \\
cmd\_inject     &  1 &  1 &  1 &  1.0 \\
reverse\_shell  &  1 &  1 &  1 &  1.0 \\
\bottomrule
\end{tabular}
\end{table}

The threshold-based rules (Nmap, SSH brute-force, ICMP flood)
produce multiple alerts per burst once the threshold is crossed,
which is the intended behaviour: each subsequent packet that keeps
the sliding window above the count limit produces another alert.
The content-based rules (SQLi, command injection, reverse shell)
produce one alert per matching packet, again as expected, because
the burst contains exactly one packet matching the regex or
content predicate.

The interesting figure is the gap between minimum and maximum
alerts for the threshold rules. For Nmap, where the burst is
between 25 and 40 packets and the threshold is 8 in 5 seconds, the
number of alerts equals \(N - 7\) for a burst of size \(N\). The
random burst size produces the spread.

\subsection{Hybrid Detection}
\label{sec:val-hybrid}

After training and activating a Random Forest model
(\code{POST /api/ml/train} with default parameters), I re-ran the
same attack script with ML scoring enabled. Table~\ref{tab:val-hybrid}
breaks down the detection methods for each attack kind.

\begin{table}[h]
\centering
\caption{Detection method breakdown with Random Forest activated.}
\label{tab:val-hybrid}
\small
\begin{tabular}{l r r r r}
\toprule
Attack kind     & Total & signature & ml only & hybrid \\
\midrule
nmap            & 137 &  47 & 0 &  90 \\
ssh\_brute      &  55 &  18 & 0 &  37 \\
icmp\_flood     & 273 & 109 & 0 & 164 \\
sqli            &   5 &   2 & 0 &   3 \\
cmd\_inject     &   5 &   2 & 0 &   3 \\
reverse\_shell  &   5 &   2 & 0 &   3 \\
\bottomrule
\end{tabular}
\end{table}

The \code{ml only} column is zero across the board because the
signature engine catches the same patterns the model was trained
on; for an ml-only alert the signature engine has to miss while
the model still finds the packet anomalous. The synthetic
catalogue is not designed to produce that scenario --- every
attack in it has a corresponding signature. The hybrid column,
however, shows the value of the cooperation: most alerts are
upgraded from signature to hybrid because the model agrees with
the rule.

\section{Performance Measurements}
\label{sec:val-performance}

A few numbers worth recording, all collected while the
synthetic catalogue was being driven through the dispatcher.

\subsection{Per-Packet Cost}
\label{sec:val-perpacket}

I instrumented the dispatcher's \code{handle()} method with
\code{time.perf\_counter()} probes and recorded the per-packet
cost across 5000 packets generated by repeated invocations of
the synthetic ICMP flood. The cost broke down as follows:

\begin{table}[h]
\centering
\caption{Per-packet processing cost (median of 5000 samples).}
\label{tab:val-perpacket}
\small
\begin{tabular}{l r r}
\toprule
Stage                                & Median (µs) & 95th percentile (µs) \\
\midrule
Ring buffer + counters               &   2 &    5 \\
Traffic bus publish                  &  18 &   32 \\
Signature evaluation (10 rules)      &  42 &   78 \\
ML scoring (Random Forest, active)   & 380 &  610 \\
Persistence + broadcast (per alert)  & 2900 & 4400 \\
\bottomrule
\end{tabular}
\end{table}

The persistence cost is the dominant one and explains why the
synthetic benchmarks above never report sub-millisecond
end-to-end latency: a single alert involves three database
round-trips (the alert insert, the file write, the PcapRecord
insert) before the WebSocket broadcast can happen. For a
thesis-scale workload this is acceptable. For a production
deployment one would batch the writes or use a faster storage
backend.

\subsection{End-to-End Latency}
\label{sec:val-latency}

The user-visible latency that matters is the time from packet
arrival to alert appearing in the browser. I instrumented this
end-to-end by injecting synthetic packets timestamped at
generation and reading the timestamp on the WebSocket message
received by a small test client. Over 200 samples in a quiet
LAN environment, the distribution was:

\begin{itemize}
  \item Median: 158\,ms.
  \item 95th percentile: 236\,ms.
  \item Maximum (excluding one outlier of 1.2\,s during a venv
    package install in the background): 412\,ms.
\end{itemize}

The 250\,ms figure repeatedly cited elsewhere in this thesis is
the 95th-percentile measurement.

\subsection{Memory Footprint}
\label{sec:val-memory}

The Python backend process steady-state memory consumption,
measured with the Windows Task Manager after a few minutes of
synthetic traffic and one trained Random Forest in memory, was
182\,MB resident. The Vite development server was a separate
Node.js process at 290\,MB resident; the production build does
not need a running Node process and would idle close to zero
once served by a static file server. Total footprint of the
running system, in the configuration used for development, was
just under half a gigabyte of RAM.

\subsection{Latency of Hybrid Categorisation}
\label{sec:val-hybrid-latency}

The hybrid categorisation adds the ML scoring path to the
critical section. With Random Forest active, the median
end-to-end latency increased to 191\,ms (from 158\,ms without
ML), the 95th percentile to 281\,ms. The increment is dominated
by the ML inference itself (median 380\,µs from the
per-packet table), amortised over the asynchronous
persistence steps. For Isolation Forest the equivalent figures
were 175\,ms median and 254\,ms p95 --- IsolationForest is
slightly faster than Random Forest because its
\code{predict()} only walks a single set of trees instead of
querying probability estimates.

\section{Discussion}
\label{sec:val-discussion}

The validation results support three claims I would defend
against a sceptical reviewer.

\textbf{The detection pipeline produces correct, repeatable
results on every attack class in the catalogue.}
Table~\ref{tab:val-synthetic} shows non-zero alert counts for
all six attack kinds, and the run-to-run variance comes from the
randomised burst sizes inside the simulator, not from
detection-engine flakiness. The threshold rules fire the
expected number of times for a burst of size $N$ (namely $N-c+1$
where $c$ is the threshold count); the content rules fire once
per matching packet. These are the algorithmically correct
behaviours.

\textbf{The hybrid categorisation works as designed.}
Table~\ref{tab:val-hybrid} shows that the Random Forest model
agrees with the signature engine on approximately two-thirds of
attack packets in the synthetic test. The remaining one-third
fall in the signature-only column because the ML model's
threshold (0.8 by default) is conservative; lowering the
threshold would convert more signature alerts into hybrid ones
at the cost of false-positive risk.

\textbf{The performance budget is sufficient for the operational
scenarios this thesis targets.} A laptop processing several
thousand packets per second through the dispatcher at
95th-percentile latency under 300\,ms is not a production-grade
deployment, but it is comfortably above the throughput one
would see on a small office network or a home lab. Scaling
beyond that point is a question of swapping the storage backend
(SQLite is the obvious bottleneck) and is not a fundamental
limitation of the architecture.

The most honest limitation is one that the synthetic
methodology cannot resolve on its own: the validation does not
exercise the machine-learning detector against attacks it was
\emph{not} trained on. A more rigorous evaluation would replay
the CICIDS2017 dataset through the dispatcher and report the
ml-only alert rate on attack categories the synthetic training
data did not contain, plus drive real attack tooling (e.g.\
\code{nmap}, \code{hydra}) from an attacker host against the
live-capture path to demonstrate parity with the synthetic
results. Both of those scenarios are on the future-work list in
Chapter~\ref{chap:conclusions} and are the natural next steps
for extending this thesis past its current scope.
